\documentclass[preprint,12pt,authoryear]{elsarticle}

\usepackage[T1]{fontenc}
\usepackage[utf8]{inputenc}
\usepackage{amsmath}
\usepackage{amssymb}
\usepackage{graphicx}
\usepackage{booktabs}
\usepackage{siunitx}
\usepackage{xspace}
\usepackage{hyperref}
\usepackage{url}

\journal{Chemical Engineering Science}

\emergencystretch=2em
\makeatletter
\setlength{\@fptop}{0pt}
\makeatother

\newcommand{\LiSi}{Li$_4$SiO$_4$\xspace}
\newcommand{\dfdp}{$d_f/d_p$\xspace}
\newcommand{\ug}{$u_g$\xspace}
\newcommand{\FdW}{$F_d/W$\xspace}
\newcommand{\BTC}{BTC\xspace}
\newcommand{\Ib}{$I_b$\xspace}

\begin{document}

\begin{frontmatter}

\title{\texorpdfstring{Separating breakthrough from pore-network clogging indicators in gas-driven fine-debris transport through \LiSi pebble beds}{Separating breakthrough from pore-network clogging indicators in gas-driven fine-debris transport through Li4SiO4 pebble beds}}

\author[aff1,aff2]{Jian Wang\corref{cor1}}
\ead{wjfttt@mail.ustc.edu.cn}
\author[aff2]{Mingzhun Lei}
\author[aff1]{Haoxi Wang}
\author[aff1]{Zhiyuan Liu}
\author[aff1]{Haishun Deng}
\author[aff1,aff2]{Wei Wen}
\author[aff1]{Gang Shen}

\cortext[cor1]{Corresponding author.}

\address[aff1]{Anhui University of Science and Technology, Huainan, Anhui, China}
\address[aff2]{Institute of Plasma Physics, Hefei Institutes of Physical Science, Chinese Academy of Sciences, Hefei, Anhui, China}

\begin{abstract}
Fine debris in gas-purged ceramic breeder beds can contaminate downstream regions before a bed develops resolved hydraulic degradation, yet outlet breakthrough is often interpreted as part of a single clogging sequence. This can overstate hydraulic risk when a mobile debris tail reaches connected pathways without observed loss of purge access. Here we examine conditional transport--filtering separation with particle-resolved DEM debris tracking and DEM-derived pore reconstruction in a physically stiff \LiSi bed containing 10,000 \SI{1}{mm} pebbles. The reconstructed baseline has a porosity of 0.4114, a largest connected void fraction of 0.9921 and a box-counting fractal dimension of 2.694. Within the limited parameter window considered here, increasing the Stokes drag-to-weight ratio shifts the debris population downstream and raises the final \BTC from zero to 0.0820, whereas no measurable degradation of the outlet-connected pore skeleton is resolved under the present DEM assumptions. A high-inventory localized-release case shows a second conditional separation: a rare front approaches the outlet ($x_{\max}/L=0.9946$ at \SI{10.00}{ms}) while the population bulk remains near the source ($x_{99}/L=0.4202$) and no particle crosses the outlet in the simulated window. Increasing the injected debris count from 3000 to 10000 raises the maximum local sub-voxel blockage from $1.1\times10^{-5}$ to $3.5\times10^{-5}$, yet the exploratory pressure-free index remains below $1.8\times10^{-5}$ because connectivity loss is not resolved in the present reconstruction. Returned cropped-domain OpenFOAM checks give a maximum 0.21\% pressure-drop increase across the same peak-blockage handoff cases, which bounds but does not validate the hydraulic response and does not close the gas--solid feedback loop. These results are consistent with a pre-clogging migration/filtering regime in which weak breakthrough or near-outlet sparse-front formation is a transport indicator rather than, by itself, evidence of bed-scale pore-network degradation. The work provides a topology-informed workflow for staged clogging assessment: \BTC flags transport and contamination, local blockage flags deposition localization, connected-pore loss flags structural degradation, and pressure-drop increase is required for hydraulic confirmation. It reports no resolved hydraulic degradation in the studied window, but requires two-way coupling for closure before validated clogging criteria can be claimed.
\end{abstract}

\begin{keyword}
ceramic breeder pebble bed \sep Li4SiO4 \sep discrete element method \sep DEM-derived pore reconstruction \sep fine debris \sep structural screening index \sep breakthrough curve
\end{keyword}

\end{frontmatter}

\section{Introduction}

Gas-purged packed beds are used when a granular skeleton must remain permeable while it is exposed to fines generation, drag-driven migration and deposition. In these systems, a small mobile debris population can be retained by pore-throat filtering, deposited along tortuous pathways or transported to downstream components without necessarily producing bed-scale hydraulic degradation. Ceramic breeder pebble beds in solid-breeder blanket concepts are a high-consequence example: they must combine tritium production, mechanical stability and purge-gas transport under demanding thermal and irradiation environments \citep{Abdou2015Blanket,Giancarli2012HCPB}. In \LiSi and related breeder beds, fine debris may be generated by handling, contact damage, thermal cycling or local fracture. The same debris event can therefore appear in several non-equivalent observables: outlet breakthrough, local accumulation, pore-connectivity loss and pressure increase.

The central difficulty is that these observables are not equivalent. A breakthrough curve reports the cumulative outlet arrival of debris, but it does not say whether the bed has become hydraulically blocked. Conversely, a connected-pore metric can show that the pore skeleton is retained while a small leading tail of debris has already reached the outlet. This distinction matters for fusion-blanket risk interpretation: weak breakthrough may indicate debris mobility, whereas clogging risk requires structural and hydraulic degradation of purge pathways. Existing packed-bed gas-powder and fine-deposition studies provide important scale references for this problem \citep{Natsui2012FinesClogging,Boccardo2019PackedBedDeposition,Wang2024IJMF}, but early mechanistic studies often lack simultaneous internal imaging, particle tracking and pressure-drop measurement.

DEM calculations provide a controlled route to this separation when their assumptions are made explicit \citep{Cundall1979DEM,Kloss2012LIGGGHTS,Plimpton1995LAMMPS}. Particle coordinates can be post-processed into voxelized pore fields, from which porosity, connected components, outlet-connected void fraction, Euler-type descriptors and box-counting dimensions can be extracted. Such morphological descriptors are widely used to connect pore geometry with connectivity and transport in porous media \citep{Armstrong2019Minkowski,Dathe2005Fractal}. The resulting fields are not experimental imaging data; they are DEM-derived pore reconstructions. Their value is that they connect resolved particle trajectories to pore-space diagnostics that can later be calibrated against pressure-drop closures or CFD.

Here we use DEM-derived voxel topology to examine whether fine-debris breakthrough in a realistic \LiSi pebble bed coincides with measurable macroscopic clogging indicators in the studied window. The bed contains 10,000 physically stiff \SI{1}{mm} pebbles with no formal modulus softening, and debris is driven by one-way axial Stokes drag. Rather than using outlet arrival as a direct proxy for clogging, we measure three responses separately: \BTC, local debris accumulation and connected-pore degradation. The simulations indicate that increasing drag-to-weight ratio promotes downstream migration and weak breakthrough while outlet-connected pore degradation is not resolved by the present reconstruction. Increasing debris loading amplifies local blockage, but the current loading window remains below the pressure-free screening reference used for the exploratory \Ib. The paper therefore frames ceramic-breeder debris migration as a pre-clogging transport/filtering problem in which particle retention, pore-throat accumulation and hydraulic consequence are coupled but non-equivalent diagnostics \citep{Natsui2012FinesClogging,Boccardo2019PackedBedDeposition,Mays2005Clogging,Dressaire2017Clogging,Parvan2020Clogging}. Supplementary Table~S5 summarizes how the present evidence is positioned relative to packed-bed deposition, porous-media clogging and morphology references.

\section{Methods}

\subsection{DEM bed and debris cases}

The DEM bed used in this study was a rectangular \LiSi pebble bed with axial flow along the positive $x$ direction and gravity along $-z$. The formal 10,000-particle bed was generated in a \SI{45}{mm} by \SI{15}{mm} by \SI{28}{mm} non-periodic container and post-processed over the settled bed domain used for voxel reconstruction. The bed contained spherical breeder pebbles with diameter $d_p = \SI{1}{mm}$, density $\rho_p=\SI{2390}{kg.m^{-3}}$, Young's modulus \SI{90}{GPa} and Poisson's ratio 0.25. Normal restitution and sliding friction coefficients were 0.50 and 0.30, respectively, for all type pairs. Contacts were represented with a Hertz--Mindlin history model in LIGGGHTS-INL, and the standard hard-particle timestep was $\Delta t=\SI{1.0e-8}{s}$ with a granular timestep check. The high-inventory localized-source continuation used $\Delta t=\SI{5.0e-9}{s}$ for additional stability during dense initial relaxation.

The bed was compacted to a random packed state and then used as the common skeleton for debris-transport calculations. Formal transport cases used the compacted skeleton as a fixed bed, implemented by freezing the type-1 large pebbles while integrating the type-2 debris particles, so that changes in breakthrough and deposition could be attributed to debris size, gas velocity and loading rather than rearrangement of the large pebbles. This choice also matches the intended mechanism window: early debris migration through an already packed breeder bed before enough fines have accumulated to restructure the pebble skeleton. It should therefore be read as a pre-clogging transport approximation, not as a model of late-stage mechanical bed deformation or debris-induced large-pebble rearrangement.

Debris particles were represented as type-2 spheres with diameter $d_f = (d_f/d_p)d_p$. The simplified purge-gas forcing used a one-way Stokes drag law,
\begin{equation}
F_{\mathrm{drag},x} = 3 \pi \mu_g d_f (u_g - v_x),
\end{equation}
applied along the axial direction. This one-way forcing provides a physically interpretable migration scale based on debris size and gas velocity without introducing an artificial force multiplier. The model intentionally neglects pore-scale gas redistribution around individual debris particles, so its role is to order migration states by a transparent Stokes-scale drive rather than to reproduce a resolved two-way gas--solid flow field. The drag-to-weight ratio was computed as
\begin{equation}
\frac{F_d}{W} = \frac{18\mu_g u_g}{\rho_d g d_f^2},
\end{equation}
and used as a compact measure of the migration drive. The corresponding debris diameters, nominal particle Reynolds numbers based on the superficial gas speed, and recomputed Stokes drag-to-weight ratios are reported in Supplementary Table~S18 as a dimensional audit of the one-way forcing model; this audit does not replace resolved pore-scale gas-flow calculation.

\subsection{Breakthrough and local blockage metrics}

Breakthrough was quantified by counting unique debris particles that crossed the downstream outlet plane. The breakthrough curve was defined as
\begin{equation}
\mathrm{BTC}(t)=\frac{N_{\mathrm{out}}(t)}{N_{\mathrm{in}}},
\end{equation}
where $N_{\mathrm{out}}(t)$ is the cumulative outlet count and $N_{\mathrm{in}}$ is the injected debris count. Local blockage was estimated by binning the debris volume along the axial direction and normalizing it by the local pore volume. This produces an axial profile $B(x)$ and a maximum local sub-voxel blockage ratio.

\subsection{DEM-derived voxel reconstruction and topology metrics}

The large-pebble coordinates were voxelized into a three-dimensional array in which pore voxels and solid voxels were separated. The baseline reconstruction used a voxel size of \SI{0.25}{mm} and a domain of \SI{45}{mm} by \SI{15}{mm} by \SI{13.25}{mm}, giving a voxel shape of $180 \times 60 \times 53$. The pore phase was treated as a binary void mask for connected-component analysis. We report total porosity, largest connected void fraction, outlet-connected void fraction, Euler number, normalized topological charge and box-counting fractal dimension.

To check whether the baseline topology was controlled by the chosen voxel size, the compacted-bed dump was independently re-voxelized at voxel sizes from \SI{0.20}{mm} to \SI{0.50}{mm}. The \SI{0.25}{mm} reconstruction matched the baseline pore metrics, while the outlet-connected void fraction remained above 0.95 for voxel sizes up to \SI{0.40}{mm}. At \SI{0.50}{mm}, the outlet-connected fraction decreased to 0.743, indicating that this coarser resolution under-resolves pore-throat connectivity. Thus, topology statements are resolution-bounded: the \SI{0.25}{mm} grid corresponds to $0.25d_p$ and cannot resolve connectivity changes smaller than approximately one voxel. A separate coarse-graining stress test was also performed on the native voxel field to quantify thresholding sensitivity.

The topology metrics are used as structural diagnostics of the DEM-derived reconstruction. The outlet-connected void fraction measures whether a pore pathway remains connected to the outlet face in the reconstructed skeleton. The topological charge is a normalized Euler-number descriptor used here as a compact indicator of pore-network complexity, following the broader use of Minkowski-functional and Euler-characteristic descriptors in porous-media morphology \citep{Armstrong2019Minkowski}. The box-counting dimension is obtained from the slope of $\log N$ against $\log(1/s)$ over sampled box sizes $s$, consistent with fractal descriptions of solid and pore phases in porous media \citep{Dathe2005Fractal}.

\subsection{Pressure-free structural screening index}

A validated clogging criterion should combine local blockage, connected-pore loss and pressure increase because particle retention can alter both local pore access and macroscopic hydraulic conductance \citep{Mays2005Clogging,Parvan2020Clogging}. In the present work, however, pressure is not solved during the DEM transport calculation and the returned OpenFOAM cases are limited cropped-domain checks. We therefore use \Ib only as a purely diagnostic, pressure-free structural indicator,
\begin{equation}
I_b = \frac{1}{2}B_{\max} + \frac{1}{2}C_{\mathrm{loss}},
\end{equation}
where $B_{\max}$ is the maximum local sub-voxel blockage ratio and $C_{\mathrm{loss}} = 1 - C_t/C_0$ is the outlet-connectivity loss relative to a reference state. The equal weights are a heuristic normalization that gives the local-deposition and connectivity-loss channels comparable bookkeeping roles; they are not fitted coefficients, not a permeability model and not a safety classification. \Ib is not a monotonic failure metric, because local blockage, connectivity loss and hydraulic resistance need not evolve on the same coordinate. It is also not comparable across bed geometries, voxel sizes or release modes without recalibration of the reference state and weighting. The value $I_b = 0.3$ is retained only as a default software screening reference for flagging cases that would deserve additional hydraulic analysis. It is not used here as a critical threshold, a validated clogging criterion or a hydraulic safety limit. \Ib cannot be used as a hydraulic failure criterion. Pressure-calibrated classification would require pressure-gradient measurements or coupled flow calculations, for example through packed-bed closures or resolved CFD/LBM \citep{Carman1937Permeability,Ergun1952PackedColumns,Wakao1982PackedBeds}.

Hydraulic consequence was evaluated in three bounded levels. First, the loading cases were post-processed with the conventional Ergun relation using the baseline voxel porosity $\epsilon_0 = 0.4114$ and a local effective porosity $\epsilon_t(x)=\epsilon_0[1-B(x)]$. Second, scalar Darcy--Laplace solves on the connected pore voxels mapped the axial blockage profile and a three-dimensional hydraulic-resistance multiplier field to conductance changes; these scalar solves are not Navier--Stokes CFD or LBM calculations. Third, three cropped-domain OpenFOAM handoff cases consumed the same resistance fields through a Darcy--Forchheimer porous-force term for the peak-blockage regions. This hierarchy is used to ask a limited question: whether the measured local debris accumulation is large enough to produce a detectable hydraulic perturbation in the present numerical window. The Ergun and scalar voxel-network results are hydraulic proxies, not CFD pressure results, and the returned OpenFOAM cases are numerical checks for cropped domains; they are not experimental pressure validation, not blanket-scale flow simulations and not used to calibrate \Ib. The pressure/flow evidence levels are separated explicitly in Supplementary Table~S19 so that structural screening, closure proxies, scalar pore-network checks and cropped-domain OpenFOAM checks are not collapsed into a validation-grade flow claim. The source data, sensitivity tests, OpenFOAM model/mesh rendering and boundary flags are provided in Supplementary Figs.~S3, S12, S15, S18, S22, S24 and S26 and Supplementary Tables~S6, S7, S10, S11, S15 and S19.

To make the mechanism interpretation auditable, the source-traced rows in the derived dimensionless map were also assigned rule-based state labels. A row was marked as breakthrough-active when final \BTC $\geq 0.01$, locally deposited when $B_{\max}\geq10^{-5}$, connectivity-loss active when $C_{\mathrm{loss}}\geq0.01$, and sparse-front active when $x_{\max}/L\geq0.95$ together with $x_{\max}-x_{99}\geq0.25L$. Localized-source rows were additionally marked as source-retained when more than half of the debris remained in the source slab. These thresholds are bookkeeping criteria used to make the case interpretation reproducible; they are not fitted transition thresholds, not a universal phase diagram and not pressure-calibrated \Ib evidence.

\section{Results}

\subsection{A connected pore skeleton provides the structural reference}

The first requirement for separating breakthrough from clogging is a pore-network reference that is independent of the outlet signal. Figure~\ref{fig:workflow} summarizes this evidence logic as three diagnostics with different physical meanings: \BTC quantifies outlet arrival, $B_{\max}$ quantifies local debris loading, and $C_{\mathrm{loss}}$ quantifies outlet-connected skeleton loss. In the present window, the largest observed values are \BTC = 0.082 and $B_{\max}=3.55\times10^{-5}$, while $C_{\mathrm{loss}}$ is not resolved above the reconstruction threshold. We therefore converted the DEM particle coordinates into a binary voxel representation of the solid and void phases before interpreting debris motion. The baseline bed contained 10,000 \LiSi pebbles in a rectangular axial-flow domain with a voxel size of \SI{0.25}{mm}. The resulting voxel array had a shape of $180 \times 60 \times 53$ and a porosity of 0.4114 (Fig.~\ref{fig:baseline}).

The baseline pore space was dominated by a connected void skeleton. The largest connected void component occupied 0.9921 of the pore phase, and the outlet-connected void fraction was also 0.9921. This indicates that, before debris loading, the reconstructed pore network provides continuous pathways across the bed. The Euler number was -18225 and the corresponding normalized topological charge was -0.03184, consistent with a multiply connected pore network rather than isolated pores. A box-counting analysis gave a fractal dimension of 2.694 with a fit $R^2$ of 0.996, showing that the voxelized pore geometry has a reproducible scale-dependent structure over the sampled box sizes.

These baseline quantities define the structural reference against which all subsequent debris cases are judged: breakthrough can be counted at the outlet, but clogging requires a measurable change in the connected pore skeleton. The DEM-derived voxel reconstructions are not experimental computed-tomography data and should be interpreted only as reconstruction metrics computed from the simulated particle coordinates.

\begin{figure*}[!htbp]
\centering
\includegraphics[width=\textwidth]{../figures/paper2_fig1_voxel_topology_framework.pdf}
\caption{DEM-derived voxel-topology diagnostic map for separating breakthrough from clogging in gas-purged pebble beds. Outlet arrival, maximum local blockage and outlet-connectivity loss are plotted as distinct response channels before the baseline porosity, connected-void metrics, drive response and inventory response are interpreted.}
\label{fig:workflow}
\end{figure*}

\begin{figure*}[!htbp]
\centering
\includegraphics[width=\textwidth]{../figures/paper2_fig2_baseline_voxel_topology.pdf}
\caption{DEM-derived structural reference for the 10,000-particle \LiSi bed. The porosity profile, transverse-averaged void map, box-counting fit and topology metrics define the connected baseline used to detect debris-induced connectivity loss.}
\label{fig:baseline}
\end{figure*}

\subsection{Drive-controlled penetration occurs without resolved connectivity loss}

Figure~\ref{fig:representative} isolates the drive axis. Three representative final states were compared using the Stokes drag-to-weight ratio \FdW, which combines debris size and purge-gas velocity into a single forcing scale. As \FdW increased from 24.08 to 59.45 and 112.86, final \BTC increased from 0 to 0.0183 and 0.0820, while the mean debris position shifted from $x/L=0.208$ to 0.358 and 0.503. The blockage centroid followed the same downstream displacement, moving from $x/L=0.208$ to 0.357 and 0.500.

The spatial profiles show how this outlet signal forms. In the low-drive state, the maximum blockage occurred at $x=\SI{8.35}{mm}$, or $x/L=0.186$, consistent with internal filtering. In the two breakthrough states, the maximum blockage shifted to $x=\SI{44.95}{mm}$, or $x/L=0.999$, indicating outlet-side enrichment through connected pathways. The figure therefore suggests a drive-controlled migration sequence within the sampled cases: stronger \FdW advances the debris and blockage fronts and produces weak breakthrough.

Formal repeat-seed calculations at $d_f/d_p = 0.0225$ and $u_g = \SI{2.0}{m.s^{-1}}$ gave final \BTC = 0.01644 $\pm$ 0.0027 and final retention = 0.9836 $\pm$ 0.0027 across the processed fixed-bed repeats (Supplementary Fig.~S17 and Supplementary Table~S14). The first-breakthrough delay was similarly bounded at \SI{0.0381 \pm 0.0014}{s} after the first saved transport frame. This repeatability check is used only to bound stochastic variability in the near-threshold DEM setting and does not define a sharp transition.

The same comparison also gives the first conditional-separation check. The $x$-direction outlet-connected void fraction was 0.999528 for all three reconstructed final states, so the breakthrough increase did not coincide with outlet-connectivity loss above the present voxel-resolution threshold. The maximum local sub-voxel blockage ratio remained $O(10^{-5})$. Thus, Fig.~\ref{fig:representative} should be read as evidence for drive-controlled penetration through a reconstructed skeleton with no resolved degradation under present DEM assumptions, rather than as evidence of bed-scale pore-network clogging.

Time-resolved deposition analysis of the representative $d_f/d_p = 0.0225$, $u_g=\SI{2.0}{m.s^{-1}}$ case shows that breakthrough follows front migration rather than sudden bed-scale collapse. The largest instantaneous blockage occurred early, at \SI{7.29}{ms} and $x/L = 0.075$, before any measurable outlet signal. The first non-zero \BTC appeared later at \SI{31.59}{ms} (\SI{0.0316}{s}), when the mean debris position was $x/L = 0.517$ and the 99th percentile was still $x/L = 0.798$. The 99th percentile crossed $x/L=0.9$ at \SI{35.64}{ms}, about \SI{0.0041}{s} after the first non-zero \BTC, and by the final frame \BTC reached 0.0183 while the mean position advanced to $x/L = 0.640$ and the 99th percentile reached $x/L = 0.997$. A pre-breakthrough fit gave a mean-cloud speed of \SI{0.78}{m.s^{-1}} and a 99th-percentile front speed of \SI{1.13}{m.s^{-1}} with $R^2=0.998$. Thus, outlet arrival is better interpreted as a delayed leading-tail event through retained pathways than as abrupt bed-scale pore-network collapse.

Together, the representative and time-resolved evidence support a specific mechanism: stronger drive promotes penetration and weak breakthrough, while topology-level degradation remains absent in the present window.

\begin{figure*}[!htbp]
\centering
\includegraphics[width=\textwidth]{../figures/paper2_fig3_representative_debris_blockage.pdf}
\caption{Drive-controlled penetration without resolved pore-connectivity loss in the studied window. Increasing \FdW shifts debris and blockage fronts downstream and raises final \BTC, whereas the reconstructed outlet-connected void fraction remains 0.999528.}
\label{fig:representative}
\end{figure*}

\subsection{Localized high-inventory release produces sparse-front saturation}

Explicit internal-release calculations test whether a localized source changes the response from drive-controlled penetration to source-zone retention. Debris particles with $d_f/d_p=0.1$ were initialized from pore-space samples in internal source slabs, while the large-pebble skeleton remained fixed and the same Stokes drag law was applied. The target-time downstream-source case retained 89.61\% of debris in the source slab at \SI{10.05}{ms}, moved 10.09\% downstream and produced no outlet release. The upstream-source case was extended to \SI{40}{ms}: 57.93\% remained in the source slab, 42.03\% moved downstream, and the outlet count still remained zero. These cases support a retention-dominated source response with slow downstream leakage, but they do not define a complete source-position scaling law.

The high-inventory upstream-source case gives a stricter population-level test because it used 30000 debris particles in the same fixed skeleton. At the completed target-time frame (1.99M steps, \SI{10.00}{ms}), the maximum debris coordinate was $x_{\max}/L=0.9946$, but the population bulk remained near the source: the median, 90th percentile and 99th percentile were $x/L=0.3455$, 0.4002 and 0.4202, respectively (Fig.~\ref{fig:localized908}). The source-slab retention fraction was 0.8945 and the downstream fraction was 0.1012. Only 14 particles lay beyond $x/L=0.80$, only two particles lay beyond $x/L=0.95$, and no particle crossed the outlet plane.

This population split is the key mechanism signal. The maximum coordinate alone would suggest that a front has nearly reached the outlet, but the quantile and tail-count evidence show a rare leading subset rather than a coherent breakthrough front. The maximum front peaked at 1.62M steps with $x_{\max}/L=0.9988$ and remained below that peak at the latest frame, indicating near-outlet sparse-front saturation rather than coherent outlet breakthrough. This result is therefore used as bounded mechanism evidence for retained-bulk/sparse-tail conditional separation, not as a critical-clogging threshold.

A cross-case mechanism-axis comparison further separates source release from sparse-front formation (Fig.~\ref{fig:localized_time}). The 906 upstream case releases 42.07\% of debris over \SI{40}{ms}, whereas the 907 downstream and 908 high-inventory cases release about 10.39\% and 10.50\% over their \SI{10}{ms} windows. No localized case shows outlet crossing. The largest final front-bulk separation occurs in 908, $x_{\max}-x_{99}=0.575L$, despite a released fraction comparable to the 907 case. Thus, the near-outlet coordinate in 908 is not a consequence of bulk washout; it is a sparse-tail event superposed on a retained source-zone population.
Figure~\ref{fig:sparse_front_diagnostics} is therefore used only as an intra-case diagnostic for 908. It shows the quantile fan and logarithmic tail fractions that distinguish a rare leading subset from population-scale release.

\begin{figure*}[!htbp]
\centering
\includegraphics[width=0.96\textwidth]{../figures/paper2_908_spatial_distribution_evidence.pdf}
\caption{Population splitting in the high-inventory localized-release case. At 1.98M steps, a rare leading front remains near the outlet side, but the debris population remains concentrated near the source: $x_{99}/L=0.4201$, $x_{\max}/L=0.9947$, and no particle crosses the outlet plane.}
\label{fig:localized908}
\end{figure*}

\begin{figure*}[!htbp]
\centering
\includegraphics[width=0.94\textwidth]{../figures/paper2_fig7_localized_time_dynamics.pdf}
\caption{Time-resolved localized-release dynamics. Source retention, downstream release, front quantiles and the released-fraction/front-gap phase plot show that downstream leakage and sparse-front formation evolve as separate responses. The extended 906 run reaches \SI{40}{ms} without outlet crossing, while the near-target 908 high-inventory case retains most particles near the source even as a rare front approaches the outlet.}
\label{fig:localized_time}
\end{figure*}

\subsection{Debris inventory raises local blockage but not skeleton-scale clogging}

Figure~\ref{fig:loading} isolates the loading axis at fixed $d_f/d_p=0.0225$ and $u_g=\SI{2.0}{m.s^{-1}}$. Three injected debris counts were compared: 3000, 6000 and 10000. The time-resolved \BTC curves showed late outlet arrival in all cases, with final outlet counts of 55, 85 and 156 particles and final \BTC values of 0.0183, 0.0142 and 0.0156, respectively. Because these are single-seed loading states with different injection histories, the final \BTC values are shown as discrete terminal responses rather than a fitted loading law.

The loading response is clearer in the deposition metric than in final \BTC. The maximum sub-voxel blockage ratio increased from $1.1\times10^{-5}$ for 3000 injected debris particles to $2.1\times10^{-5}$ and $3.5\times10^{-5}$ for 6000 and 10000 particles. The axial blockage profiles show an amplified deposition peak near $x=\SI{12.4}{mm}$, while a weaker outlet-side contribution remains visible. Thus, within this three-point sensitivity window, increasing debris inventory raises local accumulation even though the absolute debris-volume contribution remains small; the result is used to rank local-loading response, not to fit a monotonic loading law or a critical inventory threshold.

The outlet-connected void fraction was 0.999528 across the loading cases within the current reconstruction precision, so the connectivity-loss term was not resolved when normalized by the 3000-particle reference. The exploratory \Ib was therefore evaluated in a pressure-free form using local blockage and connectivity loss only. The resulting values were $5.4\times10^{-6}$, $1.0\times10^{-5}$, and $1.77\times10^{-5}$, all far below the default software screening reference of 0.3. These numbers are not safety limits and not permeability-loss estimates; they show that this loading window remains a structural pre-clogging screen, not a resolved critical pore-network-degradation transition.

Hydraulic checks give the same bounded interpretation, but at a different evidence level from the structural \Ib. The baseline Ergun pressure-gradient estimate at $u_g=\SI{2.0}{m.s^{-1}}$ was $3.90\times10^{4}\,\si{Pa.m^{-1}}$. Incorporating the axial blockage profiles increased the profile-mean pressure-gradient ratio by only $3.77\times10^{-6}$ to $1.26\times10^{-5}$ across the 3000--10000 particle loading cases, with peak-local equivalent ratios up to $8.25\times10^{-5}$. The scalar voxel-network pilot gave similarly small conductance losses of $2.8\times10^{-6}$ to $9.3\times10^{-6}$ (Supplementary Tables~S6 and S7). Cropped-domain sensitivity tests increased the estimated response but kept the conductance loss in the $10^{-6}$--$10^{-4}$ range, while the three-dimensional resistance-amplification test reached at most $1.50\times10^{-2}$ at the largest amplification factor (Supplementary Figs.~S13--S15 and S22; Supplementary Tables~S8--S10 and S15).

The returned OpenFOAM calculations provide a final cropped-domain numerical check on this hydraulic scale, not a replacement for experimental pressure validation. For the 3000, 6000 and 10000 debris handoff cases, inlet-minus-outlet pressure drops were 453.99, 454.40 and 454.94 in solver units, respectively, corresponding to a maximum relative increase of 0.21\% (Supplementary Fig.~S23 and Supplementary Table~S11). The representative cropped-domain model and mesh are shown in Fig.~\ref{fig:openfoam_model_mesh_main}, with the full three-case render set provided in Supplementary Figs.~S24 and S26, to document the numerical domain and mesh structure. These pressure-flow checks therefore support the perturbative-hydraulic-response interpretation without converting the present \Ib into a pressure-calibrated safety metric.

\begin{figure*}[!htbp]
\centering
\includegraphics[width=\textwidth]{../figures/paper2_openfoam_model_mesh_main.pdf}
\caption{Cropped-domain OpenFOAM model and mesh context for the returned pressure checks. The representative N10000 peak-blockage handoff case is shown as a rendered solid-interface domain, a central mesh slice and a local mesh-detail crop. This figure documents the numerical model and mesh used for the cropped-domain pressure check; it is not experimental validation or pressure calibration for \Ib.}
\label{fig:openfoam_model_mesh_main}
\end{figure*}

Fig.~\ref{fig:loading} therefore carries a different conclusion from Fig.~\ref{fig:representative}: inventory mainly amplifies local deposition, whereas the connected skeleton does not show resolved degradation and the pressure-free screening index remains far below the exploratory reference. Reading the three Results subsections together suggests a mechanism hierarchy rather than a single progression to clogging. The drive axis primarily affects penetration and weak outlet arrival in the sampled cases; the localized-source axis primarily affects the retained-bulk/sparse-front split; and the inventory axis primarily affects the magnitude of local deposition. In all three axes, the topology response remains below the reconstruction resolution.

\begin{figure*}[!htbp]
\centering
\includegraphics[width=\textwidth]{../figures/paper2_fig4_loading_clogging_response.pdf}
\caption{Inventory-controlled local blockage without resolved skeleton-scale clogging in the studied window. At fixed $d_f/d_p = 0.0225$ and $u_g = \SI{2.0}{m.s^{-1}}$, increasing debris count strengthens local deposition, while pressure-free \Ib remains orders of magnitude below the exploratory screening reference because reconstructed outlet-connectivity degradation is not resolved. Final \BTC values are discrete single-seed terminal states, not a fitted loading law.}
\label{fig:loading}
\end{figure*}

The final two main figures have different roles. Figure~\ref{fig:synthesis} is the interpretive summary: it assigns each case family to a dominant mechanism axis and condenses the activated observables into a family-level response map. Figure~\ref{fig:response_landscape} is the traceability and assessment check behind that summary: it displays the sparse case states, the loading response, the observable fingerprint and a staged clogging-assessment framework without fitting a universal trend. The derived dimensionless map in Supplementary Fig.~S16 and Supplementary Table~S12 extends the same bookkeeping to nine source-traced rows across the drive, loading and localized-source axes; it is used to expose evidence of observable separation, not to claim a universal regime boundary. A rule-based mechanism-state classifier applied to the same rows identifies five breakthrough-with-local-deposition cases, one retention-with-local-deposition case and three source-retained sparse-front cases, with no connectivity-loss case resolved above the selected threshold (Supplementary Table~S17). A secondary cross-observable data-mining view gives the same interpretation from the source tables: log-scale drive explains most of the drive-axis \BTC variation ($R^2=0.84$ for \BTC against $\log_{10}(F_d/W)$), loading increases $B_{\max}$ nearly linearly over the three inventory states, and the largest localized front--bulk separation reaches $0.574L$ while no connectivity-loss case is resolved in the current dataset (Supplementary Fig.~S27). A separate tail-lag mining view resolves how this separation appears in time: first non-zero \BTC occurs at \SI{31.59}{ms} and remains weak at the final frame (\BTC $=0.0183$), whereas the high-inventory localized case has a final $x_{\max}-x_{99}$ separation of $0.574L$ but only $6.7\times10^{-5}$ of particles beyond $x/L=0.95$ (Supplementary Fig.~S28). Axial deposition-heterogeneity mining adds that the representative blockage profile relaxes from a localized filtration peak to a broader downstream-spread profile, with the Gini coefficient decreasing from 0.950 to 0.481 and the participation length increasing from 0.050 to 0.578 of the bed length (Supplementary Fig.~S29). The parameter-evidence coverage map in Supplementary Fig.~S25 makes this boundary explicit by showing the discrete drive-state, loading-state and localized-source evidence that supports the present mechanism interpretation without interpolating a phase map or claiming full factorial parameter coverage.

\begin{figure*}[!htbp]
\centering
\includegraphics[width=\textwidth]{../figures/paper2_fig6_mechanism_synthesis.pdf}
\caption{Interpretive mechanism synthesis for the three observable axes. Drive, localized release and debris inventory are assigned to their dominant responses: drag-to-weight ratio activates weak outlet arrival, localized release separates source retention from sparse-front advance, and debris inventory amplifies local blockage. The family-level response map condenses these responses and shows that reconstructed connectivity loss is not resolved above the present threshold in the studied window.}
\label{fig:synthesis}
\end{figure*}

\begin{figure*}[!htbp]
\centering
\includegraphics[width=\textwidth]{../figures/paper2_fig8_response_landscape.pdf}
\caption{Case-level response landscape and staged assessment interpretation supporting Fig.~\ref{fig:synthesis}. The state map places drive and loading cases in dimensionless-loading/BTC space without fitting a trend; marker area scales with local blockage. The loading scatter and column-normalized heatmap document the source-table evidence behind the mechanism interpretation. Panel d translates the mechanism evidence into a bounded assessment sequence: \BTC is a transport/contamination indicator, local blockage is a deposition-localization indicator, connected-pore loss is a structural-degradation indicator and pressure-drop increase is required for hydraulic confirmation. This sequence is an interpretation framework, not a universal transition law or pressure-calibrated clogging criterion.}
\label{fig:response_landscape}
\end{figure*}

\begin{figure*}[!htbp]
\centering
\includegraphics[width=\textwidth]{../figures/paper2_fig9_sparse_front_diagnostics.pdf}
\caption{Intra-case sparse-front diagnostics for the 908 high-inventory localized-release case. Bulk quantiles remain near the source region while the maximum coordinate approaches the outlet side, tail fractions remain small on a logarithmic scale, and the $x_{\max}-x_{99}$ separation stays large. This figure supports the rare-front interpretation for one high-inventory source case; cross-case synthesis is given separately in Fig.~\ref{fig:synthesis}.}
\label{fig:sparse_front_diagnostics}
\end{figure*}

\begin{figure*}[!htbp]
\centering
\includegraphics[width=\textwidth]{../figures/paper2_fig10_time_resolved_deposition.pdf}
\caption{Time-resolved deposition and weak breakthrough in the representative drive case. The axial blockage field, front quantiles, \BTC and local-loading proxy show an early inlet-side blockage peak, delayed first non-zero \BTC and later 99th-percentile arrival near the outlet side. The figure is a representative post-processing view and does not define a universal front-speed law or critical-clogging threshold.}
\label{fig:time_resolved_deposition}
\end{figure*}

\subsection{A clogging criterion must keep transport, topology and pressure separate}

The combined evidence leads to a specific criterion design principle: transport, topology and pressure should be kept as separate terms until their coupling is resolved. A \BTC-only threshold would classify weak outlet arrival as a potential failure signal, even though the present breakthrough states retain the outlet-connected pore skeleton. Conversely, a topology-only metric would miss the downstream debris migration and local accumulation that precede any measurable skeleton loss. A deposition-only metric would also be incomplete because the same local debris volume can have different hydraulic meaning depending on whether it blocks connected pathways. The current results therefore support a composite index in which maximum local blockage captures deposition intensity, connected-pore loss captures structural degradation, and pressure increase captures hydraulic consequence.

At this stage, only the first two structural terms are used in \Ib. The Ergun, voxel-network and returned cropped-domain OpenFOAM pressure checks provide consistency tests on the hydraulic consequence, but they do not replace experimental pressure measurement or a validated blanket-scale flow calculation. The present \Ib values should therefore be described as a structural pre-screening index rather than a pressure-calibrated hydraulic or safety criterion. Identifying a bed-scale clogging threshold will require either higher debris loading, localized internal debris generation, or coupled flow simulations that produce measurable connectivity loss and/or pressure increase.

\section{Discussion}

The central result is evidence of response separation rather than a resolved progression to clogging. Over the present parameter window, breakthrough, local deposition and pore-network degradation do not rise as a single coupled failure mode. Instead, the sampled cases separate three mechanism axes: the drag-to-weight ratio primarily affects penetration and weak outlet arrival, the localized-release axis primarily affects the retained-bulk/sparse-front split, and the debris-inventory axis primarily affects local deposition intensity. Each axis changes a measurable transport or deposition observable while outlet-connected skeleton degradation is not resolved in the DEM-derived reconstruction. The physical picture is therefore migration and filtering through an existing connected pore network in the simulated window, not observed progressive collapse of that network.

This conditional separation changes how \BTC should be interpreted for gas-purged pebble beds. A non-zero \BTC indicates that some debris has found connected pathways to the outlet, but it does not by itself indicate loss of purge-gas access. Conversely, a topology metric alone would miss moving fronts, local accumulation and outlet contamination. The diagnostic lesson is that outlet counting, deposition maps, reconstructed pore connectivity and pressure response must be treated as related but non-equivalent observables. This interpretation is consistent with porous-media clogging studies in which retention, bridge formation, pore accessibility and permeability loss occur on different structural scales \citep{Mays2005Clogging,Dressaire2017Clogging,Parvan2020Clogging}.

The time-resolved mining clarifies the mechanism behind this diagnostic separation. In the representative drive case, first outlet arrival occurs only after the 99th-percentile front has separated from the blockage centroid, and the final \BTC remains weak even when the leading quantile reaches the outlet side. In the high-inventory localized-release case, the same separation appears in a more extreme form: the maximum coordinate approaches the outlet region, but the 99th percentile and the source-retained bulk remain far upstream, leaving only a very small survival tail beyond $x/L=0.95$ (Supplementary Fig.~S28). The axial blockage field also changes character during the same window: the initial peak/mean blockage ratio decreases from 20.0 to 2.64 while the effective participation length grows from 0.050 to 0.578 of the bed length (Supplementary Fig.~S29). Thus, the mobile tail, retained bulk, local blockage centre and spatial deposition heterogeneity should be treated as different state variables. This is the mechanism-level reason why a weak \BTC signal can coexist with a nearly unchanged connected-pore skeleton and a perturbative pressure response.

The mechanism-axis view is useful because it prevents a single observable from carrying too much meaning. \FdW is a migration coordinate, not a clogging coordinate. Debris inventory is a local-loading coordinate, not a transition coordinate in the present three-case scan. Localized release exposes the strongest population split: a small sparse front can approach the outlet side while the population bulk remains source-retained (Supplementary Fig.~S19). The cross-observable synthesis and normalized response map show that source release, sparse-front advance, local blockage, pressure proxies and connectivity loss do not collapse onto one response coordinate in the present data (Supplementary Figs.~S20 and S21). The structural axis then sets the boundary condition for interpretation: if connected-pore loss remains below reconstruction resolution, weak breakthrough and localized deposition should be described as pre-clogging transport signals rather than as pore-network failure. Supplementary Table~S13 converts these statements into claim-level evidence rows with explicit boundaries, and Supplementary Table~S17 gives the corresponding rule-based state labels while preserving the boundary that these labels are not a fitted transition law.

\subsection{Engineering interpretation for clogging assessment}

For breeder-blanket design studies, the practical value of the workflow is that it separates assessment stages that are often collapsed into a single clogging label. In the present framework, \BTC is treated as a transport and contamination indicator because it records outlet arrival but does not by itself indicate loss of purge-gas access. Local blockage is treated as a deposition-localization indicator because it maps where debris volume accumulates relative to the local pore volume. Connected-pore loss is treated as a structural-degradation indicator because it asks whether the reconstructed void skeleton has lost outlet-connected pathways, subject to voxel-resolution limits. Pressure-drop increase is then required as hydraulic confirmation because only a flow response or pressure measurement can determine whether a structural or deposition change has produced an operational hydraulic consequence.

This staged interpretation is shown in Fig.~\ref{fig:response_landscape}d and is deliberately weaker than a safety or transition criterion. A case can be transport-active without being structurally degraded, locally deposited without being hydraulically confirmed, or close to a sparse-front condition without outlet release in the simulated window. The present \Ib includes local blockage and connected-pore loss only and should therefore be read as a heuristic structural screening index, not a pressure-calibrated safety criterion. The resistance-amplification and cropped-domain OpenFOAM checks show how the missing hydraulic term can be bounded before full flow coupling: conductance loss increases monotonically with amplified debris resistance, but the returned OpenFOAM pressure drop changes by a maximum of 0.21\% across the present loading cases. Supplementary Table~S19 makes this hierarchy explicit by separating structural screening, closure proxies, scalar conductance tests and cropped-domain numerical pressure checks. These calculations strengthen the pre-clogging interpretation, but they do not replace experimental pressure measurement or blanket-scale flow simulation.

\subsection{Physical closure limits and coupling requirements}

The boundaries of the present study follow directly from the modelling choices. First, the DEM transport calculation is one-way coupled: debris experiences an imposed Stokes-scale axial drag, but the gas field does not respond to debris motion, deposition or local pore restriction. The model therefore does not resolve pore-scale gas redistribution, local velocity amplification in throats, two-way momentum exchange or feedback from deposited debris to the gas field; Supplementary Table~S18 gives the dimensional force-scaling range so that this approximation is auditable rather than hidden. Second, the type-1 pebble skeleton is frozen during debris transport. This fixed-skeleton approximation is useful for isolating early debris migration through an already packed bed, but it suppresses debris-induced large-pebble rearrangement, dilation, force-chain evolution and structural collapse. It can therefore bias the response toward retained pore connectivity when debris loading becomes large enough to restructure the matrix. Third, the DEM-derived connectivity metric is resolution-limited. The baseline uses \SI{0.25}{mm} voxels, equivalent to $0.25d_p$, so interface locations, pore-throat openings and sub-voxel blockage are uncertain at approximately one voxel unless independently converged. The \SI{0.50}{mm} re-voxelization reduces the outlet-connected void fraction to 0.743. This coarse-grid change quantifies a topology-uncertainty risk: connectivity conclusions are voxel-dependent and should be read as ``no loss resolved at the selected voxel scale'', not as proof of exact pore-throat preservation. Finally, the returned OpenFOAM calculations are cropped-domain numerical checks rather than experimental pressure validation or full-device flow solutions. The loading scan contains three single-seed states, and the internal-source evidence combines an extended upstream case, a target-time downstream case and a target-time high-inventory monitoring window. The current framework is not predictive of clogging transition. It is a bounded diagnostic workflow for separating transport, deposition, topology and pressure indicators until two-way gas--solid coupling or pressure-drop experiments provide physical closure.

\subsection{Scope of applicability}

The scope of the present conclusions is the pre-clogging regime in stiff \LiSi pebble beds under the present DEM assumptions and without resolved hydraulic feedback. The results are not intended for dense clogging, mechanically jammed beds, matrix rearrangement, late-stage pore collapse or cases in which deposited debris strongly redirects the purge flow. For prediction of bed-scale transition, the framework must be extended to two-way CFD--DEM, LBM--DEM or experimentally calibrated flow calculations in which debris deposition changes the local gas field and the updated gas field feeds back on subsequent debris motion. Within that boundary, the contribution is not a universal clogging law but a mechanism-supported diagnostic framework: \BTC, local blockage, reconstructed connectivity and pressure response should be evaluated as coupled but non-equivalent observables.

\section{Conclusions}

This study develops a DEM-derived workflow for separating fine-debris mobility from pore-network clogging indicators in stiff \LiSi pebble beds. The 10,000-particle baseline reconstruction provides a connected structural reference, with porosity 0.4114, largest connected void fraction 0.9921 and fractal dimension 2.694. Against this reference, increased \FdW raises downstream penetration and final \BTC, and increased debris inventory raises local blockage. Neither response produces outlet-connected skeleton loss above the present reconstruction threshold in the simulated cases, so the exploratory pressure-free \Ib remains $O(10^{-5})$. This workflow can be used to organize pre-clogging DEM evidence into transport, deposition, structural and hydraulic-confirmation stages, but this conclusion is restricted to the studied parameter space, under present DEM assumptions and without resolved hydraulic feedback.

The internal-source cases show why this conditional separation matters mechanistically. A localized source can retain most debris near its release slab while a rare leading subset approaches the outlet side. The upstream 15000-particle case reaches \SI{40}{ms} with no outlet release, and the target-time 30000-particle case still retains 89.45\% of debris in the source slab even though two particles lie beyond $x/L=0.95$. Thus, source-zone release, sparse-front advance and outlet breakthrough should not be collapsed into a single clogging coordinate.

Hydraulic consequence remains perturbative within the present numerical window. Ergun and voxel-network pressure proxies predict small loading-induced changes, the three-dimensional resistance-amplification test keeps the maximum local conductance loss below 0.015 even at 300-fold excess resistance, and returned cropped-domain OpenFOAM checks show a maximum 0.21\% pressure-drop increase across the 3000--10000 debris handoff cases. These findings are consistent with a pre-clogging migration/filtering regime rather than a resolved critical clogging transition. However, the conclusion is limited by the frozen pebble skeleton, one-way Stokes forcing, finite voxel topology resolution and the lack of two-way gas--solid coupling or experimental pressure validation. The present framework is diagnostic rather than predictive of clogging transition, and \Ib cannot be used as a hydraulic failure criterion. Design-oriented closure requires \BTC, local blockage, connected-pore loss and pressure response from coupled simulations or measurements, rather than treating \BTC or the present pressure-free \Ib as a standalone clogging indicator.

\section*{Data availability}

The processed data and figure source tables supporting the quantitative claims are available in the public GitHub--Zenodo repository \path{https://github.com/wangjianfttt/dem-pebble-bed-debris-transport}, archived at DOI \path{10.5281/zenodo.20699272}. Raw DEM dump and restart files are not included in the lightweight public deposit because of their size and are available from the corresponding author upon reasonable request.

\section*{Code availability}

Paper-specific scripts for source-data preparation and figure generation are available in the same GitHub--Zenodo repository under \path{papers/paper2_voxel_topology_clogging/scripts}. Shared DEM, voxel, topology and transport utilities are provided under \path{src/}. The archived repository version is available at DOI \path{10.5281/zenodo.20699272}.

\section*{CRediT authorship contribution statement}

Jian Wang: Conceptualization, Methodology, Software, Investigation, Formal analysis, Visualization, Writing -- original draft. Mingzhun Lei: Methodology, Project administration, Writing -- review and editing. Haoxi Wang: Methodology, Writing -- review and editing. Zhiyuan Liu: Application context, Methodology, Writing -- review and editing. Haishun Deng: Application context, Validation planning, Writing -- review and editing. Wei Wen: Data curation, Visualization, Writing -- review and editing. Gang Shen: Supervision, Project administration, Writing -- review and editing.

\section*{Declaration of competing interest}

The authors declare that they have no known competing financial interests or personal relationships that could have appeared to influence the work reported in this paper.

\section*{Acknowledgements}

This work was supported by the Anhui Provincial Natural Science Foundation (2408085QA030), the Anhui Intelligent Mine Technology and Equipment Engineering Research Center (AIMTEERC202307), the Science Foundation of the Institute of Plasma Physics, Chinese Academy of Sciences (DSJJ-2025-08), and the China Post-doctoral Science Foundation under Grant Number 2024M753266.

\bibliographystyle{elsarticle-harv}
\bibliography{../references/paper2_references}

\end{document}
